\par В билетах 12-13 зафиксируем поле $F$ и его расширение Галуа $L \supset F$.
\par \Def Пусть $F$ — поле, $L \supset F$ — его алгебраическое расширение. Элемент $\alpha \in L$ называется сепарабельным, если его минимальный многочлен $m_\alpha$ не имеет кратных
корней над $\overline{F}$. Расширение $L$ называется сепарабельным, если все его элементы сепарабельны.

\par \textbf{Теорема (основная теорема теории Галуа):} Сопоставление $K \mapsto Aut_K L$ (и обратное к нему сопоставление $H \mapsto L^H$) осуществляет биекцию между полями $K$ такими, что
$F \subset K \subset L$ и подгруппами $H \leq Aut_F L$, причем для соответствующих полей $K$ и
подгрупп $H$ выполнено равенство $[Aut_F L : H]=[K:F]$. Более того, расширение $L^H \supset F$ нормально $\Leftrightarrow H \trianglelefteq Aut_F L$.
\par \textbf{Доказательство биективности:} \Proof Проверим, что сопоставления $K \mapsto Aut_K L$ и $H \mapsto L^H$ взаимно обратны:
\begin{itemize}
    \item Пусть $K \mapsto Aut_K L \mapsto L^{Aut_K L}$. Расширение $L \supset K$ нормально в силу нормальности расширения $L \supset F$ (билет 58 на уд), а также сепарабельно и конечно, поэтому $L^{Aut_K L}=K$.
    
    \item Пусть $H \mapsto L^H \mapsto Aut_{L^H} L$. Очевидно, $H \subset Aut_{L^H} L$. Предположим, что $H \subsetneq Aut_{L^H} L$. По теореме о примитивном элементе, $L=L^H(\gamma)$ для некоторого $\gamma \in L$. Рассмотрим многочлен $$f(x):= \prod_{h \in H} (x-h(\gamma)) \in L[x]$$. Поскольку многочлен $f$ сохраняется под действием любого автоморфизма из $H$, то все его коэффициенты лежат в $L^H$. Но $f(\gamma)=0$, поэтому минимальный многочлен $m_\gamma \in L^H[x]$ делит $f$. Тогда, поскольку $L \supset L^H$ является расширением Галуа, выполнено $\deg m_\gamma \leq \deg f = |H| < |Aut_{L^H} L|=[L:L^H]=\deg m_\gamma$ — противоречие. 
\end{itemize}
\par Значит сопоставление является биекцией \EndProof